% GENERATED FILE. Edit canonical lessons, then run npm run generate:books.
% canonical-source-hash: fnv1a64:914284fce79b9cd4
% canonical-lessons: LA-C193-tympanum, LA-C193-talus, LA-C193-tessera, LA-C193-armilla, LA-C193-monile

\chapter{Drum, Knucklebone, Gaming die, Bracelet, Necklace}
\label{ch:la-drum-knucklebone-gaming-die-bracelet-necklace}
\hlchaptermodality{193}

\begin{chapteropening}
\textbf{By the end of this chapter:} \emph{I can say a drum, a knucklebone, a gaming die, a bracelet and a necklace.}

Complete the last lesson of chapter 193: I can say a drum, a knucklebone, a gaming die, a bracelet and a necklace.
\end{chapteropening}

\section[tympanum, tympanī]{tympanum, tympanī — a drum}

\label{lesson:LA-C193-tympanum}



\begin{quote}
Before the new one: say the Latin for an owl, then the Latin for a peacock.
\end{quote}

\subsection*{You'll want to know: tympanum, tympanī}
\textbf{tympanum, tympanī} — \textquotedblleft{}a drum\textquotedblright{}.

\begin{cousinweb}
\textbf{tympanum} was borrowed from Greek \emph{tympanon}, \textquotedblleft{}a drum\textquotedblright{} (Lewis \& Short), from \emph{typtō}, \textquotedblleft{}to strike, to beat\textquotedblright{} (EtymDB 2.1).
\end{cousinweb}

\begin{grammarlens}[title={What you've built}]
One more word for this chapter.
\end{grammarlens}

\subsection*{Your turn}
\begin{itemize}
\raggedright
  \item \emph{Say it:} \emph{tympanum, tympanī}
  \item \emph{Say it:} \emph{tympanum, tympanī}, once more
  \item \emph{Say it:} say \emph{tympanum, tympanī}
  \item \emph{Recall:} say the Latin for an owl, then the Latin for a peacock, then say \emph{tympanum, tympanī} again
\end{itemize}

\begin{tcolorbox}[breakable,colback=teal!4,colframe=teal!35!black,title={Before you move on}]
What does tympanum, tympanī mean? (\textquotedblleft{}A drum\textquotedblright{}.) Say it once more.
\end{tcolorbox}
\section[tālus, tālī]{tālus, tālī — a knucklebone; an ankle}

\label{lesson:LA-C193-talus}



\begin{quote}
Before the new one: say the Latin for a peacock, then the Latin for a drum.
\end{quote}

\subsection*{You'll want to know: tālus, tālī}
\textbf{tālus, tālī} — \textquotedblleft{}a knucklebone; an ankle\textquotedblright{}.

\begin{cousinweb}
\textbf{tālus} comes from an older \emph{*taxlus} (Lewis \& Short), as its diminutive \emph{taxillus}, \textquotedblleft{}a small die\textquotedblright{}, shows; EtymDB 2.1 gives Proto-Italic \emph{*tākslos}, \textquotedblleft{}knucklebone\textquotedblright{}, but the deeper origin is uncertain.
\end{cousinweb}

\begin{grammarlens}[title={What you've built}]
One more word for this chapter.
\end{grammarlens}

\subsection*{Your turn}
\begin{itemize}
\raggedright
  \item \emph{Say it:} \emph{tālus, tālī}
  \item \emph{Say it:} \emph{tālus, tālī}, once more
  \item \emph{Say it:} say \emph{tālus, tālī}
  \item \emph{Recall:} say the Latin for a peacock, then the Latin for a drum, then say \emph{tālus, tālī} again
\end{itemize}

\begin{tcolorbox}[breakable,colback=teal!4,colframe=teal!35!black,title={Before you move on}]
What does tālus, tālī mean? (\textquotedblleft{}A knucklebone; an ankle\textquotedblright{}.) Say it once more.
\end{tcolorbox}
\section[tessera, tesserae]{tessera, tesserae — a gaming die; a token}

\label{lesson:LA-C193-tessera}



\begin{quote}
Before the new one: say the Latin for a drum, then the Latin for a knucklebone; an ankle.
\end{quote}

\subsection*{You'll want to know: tessera, tesserae}
\textbf{tessera, tesserae} — \textquotedblleft{}a gaming die; a token\textquotedblright{}.

\begin{cousinweb}
\textbf{tessera} comes from Greek \emph{tesseres}, \textquotedblleft{}four\textquotedblright{} (Lewis \& Short): a square piece, with four corners. English \emph{tessellated}, \textquotedblleft{}made of small square tiles\textquotedblright{}, comes from its diminutive.
\end{cousinweb}

\begin{grammarlens}[title={What you've built}]
One more word for this chapter.
\end{grammarlens}

\subsection*{Your turn}
\begin{itemize}
\raggedright
  \item \emph{Say it:} \emph{tessera, tesserae}
  \item \emph{Say it:} \emph{tessera, tesserae}, once more
  \item \emph{Say it:} say \emph{tessera, tesserae}
  \item \emph{Recall:} say the Latin for a drum, then the Latin for a knucklebone; an ankle, then say \emph{tessera, tesserae} again
\end{itemize}

\begin{tcolorbox}[breakable,colback=teal!4,colframe=teal!35!black,title={Before you move on}]
What does tessera, tesserae mean? (\textquotedblleft{}A gaming die; a token\textquotedblright{}.) Say it once more.
\end{tcolorbox}
\section[armilla, armillae]{armilla, armillae — a bracelet}

\label{lesson:LA-C193-armilla}



\begin{quote}
Before the new one: say the Latin for a knucklebone; an ankle, then the Latin for a gaming die; a token.
\end{quote}

\subsection*{You'll want to know: armilla, armillae}
\textbf{armilla, armillae} — \textquotedblleft{}a bracelet\textquotedblright{}.

\begin{cousinweb}
\textbf{armilla} is built on \emph{armus}, \textquotedblleft{}the upper arm, the shoulder\textquotedblright{} (Lewis \& Short), with \emph{-illa}, a diminutive ending: a band for the arm.
\end{cousinweb}

\begin{grammarlens}[title={What you've built}]
One more word for this chapter.
\end{grammarlens}

\subsection*{Your turn}
\begin{itemize}
\raggedright
  \item \emph{Say it:} \emph{armilla, armillae}
  \item \emph{Say it:} \emph{armilla, armillae}, once more
  \item \emph{Say it:} say \emph{armilla, armillae}
  \item \emph{Recall:} say the Latin for a knucklebone; an ankle, then the Latin for a gaming die; a token, then say \emph{armilla, armillae} again
\end{itemize}

\begin{tcolorbox}[breakable,colback=teal!4,colframe=teal!35!black,title={Before you move on}]
What does armilla, armillae mean? (\textquotedblleft{}A bracelet\textquotedblright{}.) Say it once more.
\end{tcolorbox}
\section[monīle, monīlis]{monīle, monīlis — a necklace}

\label{lesson:LA-C193-monile}



\begin{quote}
Before the new one: say the Latin for a gaming die; a token, then the Latin for a bracelet.
\end{quote}

\subsection*{You'll want to know: monīle, monīlis}
\textbf{monīle, monīlis} — \textquotedblleft{}a necklace\textquotedblright{}.

\begin{cousinweb}
\textbf{monīle} is usually linked with a Proto-Indo-European word \emph{*mon-}, \textquotedblleft{}neck\textquotedblright{}, seen in Old Irish \emph{muin}, \textquotedblleft{}neck\textquotedblright{}; English \emph{mane} is a cousin.
\end{cousinweb}

\begin{grammarlens}[title={What you've built}]
One more word for this chapter.
\end{grammarlens}

\subsection*{Your turn}
\begin{itemize}
\raggedright
  \item \emph{Say it:} \emph{monīle, monīlis}
  \item \emph{Say it:} \emph{monīle, monīlis}, once more
  \item \emph{Say it:} say \emph{monīle, monīlis}
  \item \emph{Recall:} say the Latin for a gaming die; a token, then the Latin for a bracelet, then say \emph{monīle, monīlis} again
\end{itemize}

\begin{tcolorbox}[breakable,colback=teal!4,colframe=teal!35!black,title={Before you move on}]
What does monīle, monīlis mean? (\textquotedblleft{}A necklace\textquotedblright{}.) Say it once more.
\end{tcolorbox}
